\chapter{Các công trình nghiên cứu liên quan}
\label{Chapter2}

Trong giai đoạn hiện nay, rất nhiều nghiên cứu về ứng dụng học sâu đa phương thức vào lĩnh vực y tế đã phát triển rất nhanh và đạt được nhiều thành tựu nổi bật. Nhằm xác lập cơ sở cho việc thiết kế hệ thống phân lớp, phần này sẽ mô tả tổng quan các nghiên cứu về các mô hình nền tảng của các cách tiếp cận khác nhau cho bài toán phân lớp ảnh y khoa, các chiến lược bảo toàn thông tin trên ảnh độ phân giải cao, đặc biệt là cách kết hợp ngữ cảnh toàn cục với đặc trưng cục bộ, các kỹ thuật tinh chỉnh hiệu quả tham số và phương pháp phát hiện dữ liệu ngoài phân phối. Trên cơ sở đó, đối chiếu ưu điểm, hạn chế và mức độ phù hợp của các hướng tiếp cận hiện có đối với dữ liệu thực tế, chương xác định khoảng trống nghiên cứu và làm rõ cơ sở lựa chọn các thành phần trong phương pháp được đề xuất ở Chương \ref{Chapter3}.

\section{Các hướng tiếp cận chính và các mô hình nền tảng}

\subsection{Các mô hình kết hợp ngôn ngữ - ảnh xây dựng từ đầu}

Sự kết hợp giữa dữ liệu cận lâm sàng và báo cáo lâm sàng, lịch sử bệnh lý đã tạo ra những bước tiến lớn trong việc cải thiện độ chính xác của các mô hình chẩn đoán, đặc biệt là thông qua việc sử dụng văn bản để bù đắp những thiếu sót thông tin mà hình ảnh đơn thuần không thể cung cấp, hoặc giúp mô hình tập trung vào những vùng trên ảnh được xác định bởi thông tin từ nguồn văn bản. Nghiên cứu của X. Mei \cite{XMei} trong việc chẩn đoán COVID-19 là một minh chứng rõ nét: khi ảnh chụp CT ngực ở giai đoạn sớm có thể chưa biểu hiện bất thường, việc tích hợp thông tin lâm sàng thông qua mạng MLP với đặc trưng ảnh từ CNN đã giúp mô hình nhận diện chính xác 68\% số bệnh nhân dương tính mà các bác sĩ X-quang đã bỏ sót. Mô hình TieNet \cite{TieNet} đã mã hóa các báo cáo y tế và tích hợp các cơ chế chú ý đa cấp để gióng hàng các từ mang ý nghĩa quan trọng với các vùng ảnh tương ứng, cải thiện trung bình 6\% chỉ số AUC so với các mô hình chỉ dùng ảnh. Gần đây, IRENE \cite{IRENE} đưa tất cả dữ liệu về một dạng đơn vị biểu diễn chung và sử dụng một mô hình Transformer duy nhất, cho phép văn bản trực tiếp tương tác và giải thích các đơn vị biểu diễn của ảnh ở cấp độ cục bộ.

Nhìn chung, những nghiên cứu này minh chứng rằng sự đóng góp của nhiều nguồn dữ liệu, đặc biệt là văn bản đã đi từ việc chỉ là một véc-tơ thông tin ghép nối thêm, trở thành một thành phần cốt lõi có khả năng gióng hàng ngữ nghĩa và định hướng trong cơ chế chú ý. Tuy nhiên, các mô hình xây dựng từ đầu thường đòi hỏi dữ liệu ghép cặp quy mô lớn và chi phí huấn luyện cao, tạo tiền đề cho sự xuất hiện của các mô hình ngôn ngữ trực quan y khoa.

\subsection{Mô hình ngôn ngữ trực quan y khoa}

Các mô hình ngôn ngữ trực quan đánh dấu một bước tiến quan trọng trong việc ứng dụng học sâu đa phương thức vào lĩnh vực y tế. Khác với các mô hình được xây dựng từ đầu, các mô hình ngôn ngữ trực quan nói chung và các mô hình ngôn ngữ trực quan y khoa nói riêng tập trung vào hai nguồn dữ liệu đầu vào chính là ảnh và văn bản. Các mô hình ngôn ngữ trực quan y khoa kế thừa các bộ mã hóa đặc trưng độc lập cho từng phương thức, điển hình như ResNet \cite{ResNet} hoặc ViT \cite{Dosovitskiy2021ViT} cho hình ảnh và các mô hình ngôn ngữ cho văn bản. Thay vì sử dụng các mô-đun dung hợp phức tạp, các mô hình ngôn ngữ trực quan y khoa tối ưu hóa các hàm mất mát để căn chỉnh không gian đặc trưng của hai phương thức này. Nhờ biểu diễn đa phương thức mở rộng này, mô hình có khả năng hiểu sâu sắc ngữ nghĩa giữa ảnh X-quang và báo cáo lâm sàng, từ đó cải thiện đáng kể hiệu năng trên nhiều tác vụ y khoa khác nhau.

Chiến lược học tương phản là hướng tiếp cận được ứng dụng rộng rãi nhất. Nguyên lý của chiến lược này là tối ưu hóa không gian biểu diễn bằng cách thu hẹp khoảng cách giữa các cặp đặc trưng ảnh-văn bản tương đồng, đồng thời đẩy các cặp khác biệt ra xa nhau. Bảng \ref{tab:VLMPretrainingMethodPropose} tổng hợp các mô hình ngôn ngữ trực quan y khoa tiêu biểu sử dụng học tương phản.

% Bảng tổng hợp các Med-VLMs liên quan
\begin{table}[H]
\centering
\resizebox{\textwidth}{!}{
\begin{tabular}{|c|l|l|llll|}
\hline
\multirow{2}{*}{\textbf{STT}} &
  \multicolumn{1}{c|}{\multirow{2}{*}{\textbf{\begin{tabular}[c]{@{}c@{}}Tên phương \\ pháp\end{tabular}}}} &
  \multicolumn{1}{c|}{\multirow{2}{*}{\textbf{\begin{tabular}[c]{@{}c@{}}Chiến lược \\ tiền huấn \\ luyện\end{tabular}}}} &
  \multicolumn{4}{c|}{\textbf{Phương pháp}} \\ \cline{4-7}
 &
  \multicolumn{1}{c|}{} &
  \multicolumn{1}{c|}{} &
  \multicolumn{1}{c|}{\textbf{\begin{tabular}[c]{@{}c@{}}Bộ mã hoá \\ thị giác\end{tabular}}} &
  \multicolumn{1}{c|}{\textbf{\begin{tabular}[c]{@{}c@{}}Bộ mã hoá \\ văn bản\end{tabular}}} &
  \multicolumn{1}{c|}{\textbf{\begin{tabular}[c]{@{}c@{}}Chiến lược \\ kết hợp\end{tabular}}} &
  \multicolumn{1}{c|}{\textbf{Bộ giải mã}} \\ \hline
\textbf{1} &
  \textbf{\begin{tabular}[c]{@{}l@{}}ConVIRT\\ (2020)\end{tabular}} &
  \begin{tabular}[c]{@{}l@{}}Học tương\\ phản\end{tabular} &
  \multicolumn{1}{l|}{ResNet50} &
  \multicolumn{1}{l|}{BERT} &
  \multicolumn{1}{l|}{\begin{tabular}[c]{@{}l@{}}Contrastive \\ Alignment\end{tabular}} &
  N/a \\ \hline
\textbf{2} &
  \textbf{\begin{tabular}[c]{@{}l@{}}GLoRIA\\ (2021)\end{tabular}} &
  \begin{tabular}[c]{@{}l@{}}Học tương\\ phản\end{tabular} &
  \multicolumn{1}{l|}{ResNet50} &
  \multicolumn{1}{l|}{BioClinicalBERT} &
  \multicolumn{1}{l|}{\begin{tabular}[c]{@{}l@{}}Global-Local \\ Contrastive \\ Alignment\end{tabular}} &
  N/a \\ \hline
\textbf{3} &
  \textbf{\begin{tabular}[c]{@{}l@{}}MedCLIP\\ (2022)\end{tabular}} &
  \begin{tabular}[c]{@{}l@{}}Học tương\\ phản\end{tabular} &
  \multicolumn{1}{l|}{\begin{tabular}[c]{@{}l@{}}Swin Transformer, \\ ResNet50\end{tabular}} &
  \multicolumn{1}{l|}{BioClinicalBERT} &
  \multicolumn{1}{l|}{\begin{tabular}[c]{@{}l@{}}Knowledge-enhanced \\ Alignment\end{tabular}} &
  N/a \\ \hline
\textbf{4} &
  \textbf{\begin{tabular}[c]{@{}l@{}}BiomedCLIP\\ (2023)\end{tabular}} &
  \begin{tabular}[c]{@{}l@{}}Học tương\\ phản\end{tabular} &
  \multicolumn{1}{l|}{ViT-Base} &
  \multicolumn{1}{l|}{PubMedBERT} &
  \multicolumn{1}{l|}{\begin{tabular}[c]{@{}l@{}}Contrastive \\ Alignment\end{tabular}} &
  N/a \\ \hline
\end{tabular}
}
\caption{Tổng hợp các mô hình ngôn ngữ trực quan y khoa sử dụng học tương phản}
\label{tab:VLMPretrainingMethodPropose}
\end{table}

ConVIRT \cite{ConVIRT} là nghiên cứu tiên phong áp dụng học tương phản cho dữ liệu y khoa, sử dụng ResNet50 và BERT để tối ưu bộ mã hóa ảnh và văn bản, đạt độ chính xác 90,3\% trên bộ dữ liệu COVIDx. GLoRIA \cite{GLoRIA} mở rộng hướng tiếp cận này bằng cách kết hợp hai hàm mất mát toàn cục và cục bộ, cho phép căn chỉnh chi tiết giữa các vùng trong ảnh và các từ trong câu. MedCLIP \cite{MedCLIP} đề xuất cơ chế tách cặp và trích xuất tri thức cùng hàm mất mát khớp ngữ nghĩa nhằm giải quyết vấn đề âm tính giả khi các cặp dữ liệu bị ghép nhầm dù chúng có chung đặc điểm bệnh lý. BiomedCLIP \cite{BiomedCLIP} kế thừa kiến trúc CLIP \cite{CLIP} với bộ mã hóa ViT-Base và PubMedBERT, được tiền huấn luyện trên 15 triệu cặp dữ liệu từ tập PMC-15M, đạt độ chính xác phân lớp 83\%, giúp mô hình có độ phủ và khả năng tổng quát tốt với nhiều nguồn dữ liệu và tác vụ y khoa.

Ngoài học tương phản, các hướng tiếp cận khác bao gồm học có che giấu và tích hợp tri thức y khoa. Tuy nhiên, các phương pháp này thường phụ thuộc vào nguồn dữ liệu ghép cặp sẵn hoặc đồ thị tri thức chuyên biệt, hạn chế khả năng ứng dụng trên các bài toán không có sẵn nguồn tri thức ngoài. Song song với các mô hình ngôn ngữ trực quan y khoa, sự phát triển của các mô hình ngôn ngữ lớn đa phương thức như Med-Flamingo \cite{moor_2023_medflamingo}, LLaVA-Med \cite{li_2023_llavamed} hay Med-PaLM M \cite{MED-PALM-M} đã cho thấy tiềm năng trong các tác vụ hội thoại và hỗ trợ ra quyết định y khoa. Tuy nhiên, các mô hình này thường có quy mô rất lớn, chi phí huấn luyện và suy luận cao, đồng thời vẫn tồn tại nguy cơ tạo ra thông tin không chính xác. Do đó, đối với các bài toán phân lớp chuyên biệt như phân lớp ảnh X-quang xương, việc sử dụng các mô hình ngôn ngữ - thị giác chuyên dụng kết hợp tinh chỉnh hiệu quả tham số thường phù hợp và hiệu quả hơn.

\subsubsection{Lựa chọn BiomedCLIP làm mô hình nền tảng}
\label{subsubsec:BiomedCLIP_selection}
Trong số các mô hình ngôn ngữ trực quan y khoa được khảo sát, khóa luận lựa chọn BiomedCLIP \cite{BiomedCLIP} làm mô hình nền tảng dựa trên mức độ phù hợp với bài toán, thay vì chỉ dựa trên kết quả của một bộ dữ liệu riêng lẻ. Thứ nhất, BiomedCLIP được tiền huấn luyện trên PMC-15M, gồm khoảng 15 triệu cặp ảnh-văn bản thu thập từ nhiều loại dữ liệu y sinh, trong khi một số mô hình trước đó chủ yếu được xây dựng và đánh giá trên ảnh X-quang ngực. Phạm vi tiền huấn luyện rộng hơn tạo cơ sở phù hợp cho việc chuyển giao sang ảnh X-quang xương, vốn có sự đa dạng lớn về vùng giải phẫu, tư thế chụp và loại tổn thương.

Thứ hai, mô hình sử dụng kiến trúc hai bộ mã hóa độc lập gồm Vision Transformer cho ảnh và PubMedBERT cho văn bản, phù hợp với thiết kế hai nhánh của hệ thống và cho phép kế thừa trực tiếp biểu diễn của cả hai phương thức. Thứ ba, kiến trúc CLIP của mô hình hỗ trợ đánh giá không mẫu thông qua độ tương đồng ảnh-văn bản, đồng thời có thể được mở rộng bằng các kỹ thuật tinh chỉnh hiệu quả tham số và mô-đun dung hợp cho bài toán đích. Vì vậy, BiomedCLIP được sử dụng làm mô hình nền tảng khởi tạo, trong khi các ý tưởng phù hợp từ những nghiên cứu khác, chẳng hạn hàm mất mát khớp ngữ nghĩa của MedCLIP, được kế thừa để cải thiện quá trình thích nghi trên dữ liệu X-quang xương. Kết quả phân lớp không mẫu của BiomedCLIP và các mô hình nền tảng khác được trình bày trong Phần \ref{subsec:foundation_model_evaluation}.

\subsubsection{Cơ chế phân lớp không mẫu của BiomedCLIP}
\label{subsubsec:BiomedCLIP_zeroshot}

BiomedCLIP có thể thực hiện phân lớp trong thiết lập \textit{không mẫu} mà không cần huấn luyện một đầu phân lớp riêng trên dữ liệu đích. Thay vào đó, mô hình biểu diễn ảnh đầu vào và các câu nhắc mô tả lớp trong cùng một không gian đặc trưng, sau đó xác định lớp dựa trên mức độ tương đồng giữa hai phương thức.

Ảnh đầu vào $x$ trước hết được tiền xử lý theo cấu hình của BiomedCLIP, bao gồm thay đổi kích thước, cắt vùng trung tâm và chuẩn hóa, để thu được tensor ảnh có kích thước $224\times224$. Ảnh sau đó được chia thành các patch kích thước $16\times16$ và đưa qua bộ mã hóa thị giác ViT-B/16 để tạo biểu diễn ảnh.

Với mỗi lớp $k$, tên lớp được đưa vào một mẫu câu nhắc thống nhất, chẳng hạn \texttt{A radiograph showing [tên lớp].} Câu nhắc được mã hóa bởi bộ mã hóa văn bản PubMedBERT để thu được biểu diễn văn bản $t_k$. Điểm số của lớp $k$ được tính bằng độ tương đồng cosine:

\begin{equation}
s_k =
\frac{v^\top t_k}
{\lVert v\rVert_2\lVert t_k\rVert_2}.
\end{equation}

Các điểm tương đồng được chuẩn hóa bằng hàm Softmax để tạo phân phối
xác suất trên tập lớp:

\begin{equation}
p(y=k\mid x)
=
\frac{\exp(s_k/\tau)}
{\sum_{j=1}^{K}\exp(s_j/\tau)},
\end{equation}

trong đó $\tau$ là hệ số nhiệt độ. Lớp có xác suất lớn nhất được chọn làm kết quả dự đoán. Cơ chế này cho phép BiomedCLIP chuyển giao tri thức đã học sang các tập nhãn mới mà không cần cập nhật tham số. Tuy nhiên, hiệu quả phụ thuộc đáng kể vào mức độ tương đồng giữa dữ liệu tiền huấn luyện và dữ liệu đích, cách xây dựng câu nhắc, cũng như khả năng biểu diễn các dấu hiệu bệnh lý chuyên biệt của bộ mã hóa.

\section{Các kỹ thuật tinh chỉnh hiệu quả tham số}

Sự xuất hiện của các mô hình nền tảng nói chung và mô hình ngôn ngữ--thị giác y sinh nói riêng đã mang lại một nền tảng tri thức đa phương thức khổng lồ. Tuy nhiên, để triển khai hiệu quả các mô hình này vào các bài toán lâm sàng đặc thù, quá trình chuyển giao tri thức là bắt buộc.

Hướng đến yêu cầu cụ thể hơn cho bài toán sử dụng hai nguồn ảnh và văn bản lâm sàng, khóa luận tập trung vào các kỹ thuật tinh chỉnh để khai thác tối đa từ hai nguồn dữ liệu này với đặc điểm dữ liệu thực tế còn hạn chế.

Trước đây, hai chiến lược phổ biến nhất là tinh chỉnh toàn phần và dò tuyến tính. Tinh chỉnh toàn phần cập nhật lại toàn bộ tham số của mô hình gốc, mang lại khả năng thích ứng cao nhưng đòi hỏi chi phí tính toán, bộ nhớ khổng lồ, yêu cầu dữ liệu huấn luyện lớn đến rất lớn, dẫn đến hiện tượng quên biểu diễn trên các tập dữ liệu y tế nhỏ. Ngược lại, dò tuyến tính đóng băng toàn bộ mạng nơ-ron gốc và chỉ huấn luyện một lớp phân lớp duy nhất ở cuối, tiết kiệm tài nguyên nhưng hạn chế khả năng thích ứng khi sự chênh lệch phân phối giữa dữ liệu gốc và dữ liệu y tế là quá lớn.

Để giải quyết bài toán đánh đổi giữa hiệu năng và chi phí tính toán, tinh chỉnh hiệu quả tham số đã nổi lên như một giải pháp thay thế. Kỹ thuật này chỉ cập nhật hoặc thêm mới một lượng rất nhỏ tham số, giúp tiết kiệm tài nguyên và giữ vững khả năng tổng quát hóa. Dựa trên cơ chế can thiệp, các phương pháp tinh chỉnh hiệu quả tham số được chia thành ba hướng tiếp cận chính:

\subsection{Tinh chỉnh sử dụng câu nhắc}
Hướng tiếp cận này chèn thêm các tham số (vectơ) có thể học được vào không gian đầu vào của mô hình mà không thay đổi cấu trúc mạng bên trong.
\begin{itemize}
    \item \textbf{Điểm mạnh:} Tiết kiệm tham số tối đa, bảo toàn nguyên vẹn tri thức của mô hình gốc, đặc biệt hiệu quả trong các kịch bản học ít mẫu hoặc không mẫu.
    \item \textbf{Điểm yếu:} Tối ưu hóa khó khăn do không gian biểu diễn bị giới hạn ở đầu vào, nhạy cảm với cách khởi tạo, và chiếm dụng một phần độ dài của chuỗi đầu vào.
\end{itemize}
Trong lĩnh vực y tế, PLIP \cite{PLIP} kết hợp đặc trưng thị giác với câu nhắc hiệu quả, đạt 0,896 AUC chỉ với 5\% dữ liệu PanNuke. BiomedCoOp \cite{BiomedCoOp} tích hợp mô hình ngôn ngữ lớn để tự động sinh các câu nhắc ngữ nghĩa, đồng bộ hóa qua cơ chế chú ý chéo và căn chỉnh tương phản.

\subsection{Tinh chỉnh bộ điều hợp}
Phương pháp này chèn các mô-đun mạng nơ-ron siêu nhẹ, còn được gọi là bộ điều hợp vào giữa các lớp kiến trúc của mạng học sâu.
\begin{itemize}
    \item \textbf{Điểm mạnh:} Tính mô-đun hóa cao, dễ dàng thêm/bớt các bộ điều hợp cho các tác vụ khác nhau mà không can thiệp vào tham số gốc.
    \item \textbf{Điểm yếu:} Việc chèn thêm các lớp mạng vật lý làm tăng độ trễ suy luận, hạn chế khả năng triển khai trên các hệ thống đòi hỏi thời gian thực.
\end{itemize}
Điển hình, CLIPath \cite{CLIPath} sử dụng kết nối thặng dư để giảm khoảng cách miền trên dữ liệu nhỏ. FTB \cite{FTB} đặt mô-đun dung hợp giữa bộ mã hóa ảnh và văn bản để chắt lọc thông tin đa phương thức.

\subsection{Tinh chỉnh thích nghi hạng thấp}
Tinh chỉnh thích nghi hạng thấp (LoRA) giữ cố định trọng số nền và bổ sung các ma trận hạng thấp có thể huấn luyện để tham số hóa phần cập nhật trọng số. Cách tiếp cận này làm tăng một lượng nhỏ tham số huấn luyện nhưng tránh cập nhật toàn bộ mô hình.
\begin{itemize}
    \item \textbf{Điểm mạnh:} Giảm đáng kể số tham số cần cập nhật và bộ nhớ tối ưu hóa. Ngoài ra, các ma trận cập nhật có thể được hợp nhất vào trọng số nền khi cách triển khai hỗ trợ.
    \item \textbf{Điểm yếu:} Hiệu quả phụ thuộc vào vị trí chèn, hạng của ma trận và cấu hình tối ưu. Tuy nhiên, việc giảm tham số huấn luyện không đồng nghĩa với giảm chi phí lan truyền tiến của toàn bộ kiến trúc.
\end{itemize}
LCBB \cite{LCBB} áp dụng LoRA cho mô hình ngôn ngữ trực quan y khoa và đạt 93,1\% AUROC trên tập RSNA với 10\% dữ liệu gán nhãn, cho thấy khả năng thích nghi với số tham số cập nhật thấp hơn tinh chỉnh toàn phần. 

Bảng~\ref{tab:VLMTransferlearningMethodPropose} tổng hợp các phương pháp tinh chỉnh hiệu quả tham số đại diện.

% Bảng tổng hợp đại diện PEFT
\begin{table}[htbp]
\centering
\resizebox{\columnwidth}{!}{%
\begin{tabular}{|c|l|l|l|l|}
\hline
\textbf{STT} &
  \multicolumn{1}{c|}{\textbf{\begin{tabular}[c]{@{}c@{}}Tên phương \\ pháp\end{tabular}}} &
  \multicolumn{1}{c|}{\textbf{\begin{tabular}[c]{@{}c@{}}Chiến lược \\ tinh chỉnh\end{tabular}}} &
  \multicolumn{1}{c|}{\textbf{Nguyên lý}} &
  \multicolumn{1}{c|}{\textbf{Kết quả nổi bật}} \\ \hline
\textbf{1} &
  \textbf{PLIP (2024)} &
  \begin{tabular}[c]{@{}l@{}}Tinh chỉnh sử \\ dụng câu nhắc\end{tabular} &
  \begin{tabular}[c]{@{}l@{}}Kết hợp câu nhắc thị giác \\ với tinh chỉnh chọn lọc, tận \\ dụng đặc trưng hình ảnh\\ để tinh chỉnh hiệu quả\end{tabular} &
  \begin{tabular}[c]{@{}l@{}}0,896 AUC \\ (5\% PanNuke)\end{tabular} \\ \hline
\textbf{2} &
  \textbf{CLIPath (2023)} &
  \begin{tabular}[c]{@{}l@{}}Tinh chỉnh sử \\ dụng bộ điều hợp\end{tabular} &
  \begin{tabular}[c]{@{}l@{}}Sử dụng kết nối thặng dư \\ để giảm khoảng cách miền \\ trên dữ liệu mô bệnh học\end{tabular} &
  \begin{tabular}[c]{@{}l@{}}Cải thiện so với CLIP \\ trên dữ liệu nhỏ\end{tabular} \\ \hline
\textbf{3} &
  \textbf{LCBB (2024)} &
  \begin{tabular}[c]{@{}l@{}}Thích nghi hạng \\ thấp bằng LoRA\end{tabular} &
  \begin{tabular}[c]{@{}l@{}}Bổ sung ma trận hạng thấp \\ vào bộ mã hóa thị giác, giữ \\ cố định trọng số nền\end{tabular} &
  \begin{tabular}[c]{@{}l@{}}93,1\% AUROC \\ (RSNA, 10\%)\end{tabular} \\ \hline
\textbf{4} &
  \textbf{FTB (2024)} &
  \begin{tabular}[c]{@{}l@{}}Tinh chỉnh sử \\ dụng bộ điều hợp\end{tabular} &
  \begin{tabular}[c]{@{}l@{}}Mô-đun dung hợp giữa bộ \\ mã hóa ảnh và văn bản, \\ chắt lọc thông tin đa \\ phương thức\end{tabular} &
  \begin{tabular}[c]{@{}l@{}}Cải thiện trên \\ CheXpert, RSNA\end{tabular} \\ \hline
\end{tabular}%
}
\caption{Tổng hợp các phương pháp tinh chỉnh hiệu quả tham số đại diện cho các mô hình ngôn ngữ trực quan y khoa}
\label{tab:VLMTransferlearningMethodPropose}
\end{table}

\section{Các phương pháp phát hiện dữ liệu ngoài phân phối}

Phát hiện dữ liệu ngoài phân phối nhằm nhận diện các mẫu không phù hợp với phân phối quan sát trong tập huấn luyện. Trong phân tích ảnh y khoa, điểm ngoài phân phối có thể cung cấp tín hiệu cảnh báo bổ trợ khi mô hình gặp dữ liệu lạ, nhưng không tự thân bảo đảm tính an toàn hoặc thay thế quy trình đánh giá lâm sàng. Dựa trên cách tiếp cận, các phương pháp phát hiện dữ liệu ngoài phân phối hiện nay được chia thành các nhóm chính sau.

\subsection{Các phương pháp hậu xử lý dựa trên đầu ra của mô hình}
Nhóm phương pháp này không can thiệp vào quá trình huấn luyện mà tận dụng trực tiếp đầu ra (thường là các logit hoặc xác suất) của một bộ phân lớp đã được huấn luyện.

\begin{itemize}
    \item \textbf{Xác suất softmax cực đại \cite{Hendrycks2017MSP}:} Sử dụng giá trị cực đại của hàm softmax làm độ tin cậy. Dữ liệu có xác suất cực đại thấp sẽ bị đánh dấu là ngoài phân phối. Tuy nhiên, điểm yếu chí mạng của xác suất softmax cực đại là các mạng nơ-ron sâu thường mắc lỗi "tự tin thái quá", dẫn đến việc gán xác suất rất cao cho cả những mẫu ngoài phân phối hoàn toàn không liên quan.
    \item \textbf{ODIN \cite{Liang2018ODIN}:} Cải thiện xác suất softmax cực đại bằng cách sử dụng hiệu chỉnh nhiệt độ và thêm nhiễu đầu vào để gia tăng sự chênh lệch logit giữa dữ liệu trong phân phối và dữ liệu ngoài phân phối.
    \item \textbf{Điểm năng lượng \cite{Liu2020EnergyOOD}:} Chuyển đổi vectơ logit thành hàm năng lượng thay vì xác suất softmax:
    \begin{equation}
    E(x) = -T \log \sum_i \exp\left(\frac{f_i(x)}{T}\right)
    \label{eq:energy_score}
    \end{equation}
    trong đó $f_i(x)$ là logit của lớp thứ $i$ và $T$ là hệ số nhiệt độ.
    Mặc dù điểm năng lượng khai thác độ lớn tuyệt đối của các logit và thường khắc phục một phần hiện tượng quá tự tin của hàm softmax, các phương pháp hậu xử lý dựa trên đầu ra vẫn phụ thuộc mạnh vào biểu diễn và biên quyết định của bộ phân lớp đã huấn luyện. Do mô hình phân lớp đóng buộc mọi đầu vào vào một trong các lớp đã biết, các mẫu ngoài phân phối nằm xa dữ liệu huấn luyện vẫn có thể nhận điểm tin cậy cao.
\end{itemize}

\subsection{Các phương pháp dựa trên độ không chắc chắn và tái tạo}
Để tránh phụ thuộc vào đầu ra hàm softmax, một số nghiên cứu tiếp cận bài toán bằng cách đánh giá khả năng sinh hoặc sự bất định của mô hình.

\begin{itemize}
    \item \textbf{Sai số tái tạo:} Sử dụng các mô hình sinh như bộ tự mã hóa được huấn luyện chỉ trên dữ liệu trong phân phối. Khi gặp dữ liệu ngoài phân phối, sai số tái tạo sẽ tăng đột biến. Dù hiệu quả, phương pháp này gặp khó khăn khi tái tạo các ảnh y khoa có độ phân giải cao và chi tiết phức tạp.
    \item \textbf{Dropout Monte Carlo:} Kích hoạt hàm tắt ngẫu nhiên trong quá trình suy luận và thực hiện dự đoán nhiều lần để tính phương sai. Tuy nhiên, việc phải suy luận nhiều lần làm tăng chi phí tính toán, không phù hợp cho các hệ thống đòi hỏi thời gian thực.
\end{itemize}

\subsection[Các phương pháp dựa trên khoảng cách]{Các phương pháp dựa trên khoảng\\cách đặc trưng}
Nhận thấy những hạn chế của không gian logit, các nghiên cứu gần đây chuyển hướng sang thực hiện phát hiện dữ liệu ngoài phân phối trực tiếp trên không gian đặc trưng. Việc đánh giá bằng khoảng cách giúp giới hạn không gian rủi ro mở.

\begin{itemize}
    \item \textbf{Khoảng cách euclid và cosine:}
    Khoảng cách euclid đo độ lệch hình học tuyệt đối $d_E(z,c_k)=\|z-c_k\|_2$ \label{eq:euclid} và phù hợp với hình học đẳng hướng khi các chiều có cùng thang đo. Khoảng cách cosine loại bỏ ảnh hưởng của độ lớn vectơ và chỉ so sánh hướng, nhưng không mô hình hóa mức phân tán khác nhau theo từng chiều.
    \item \textbf{Khoảng cách Mahalanobis \cite{Lee2018MahalanobisOOD}:} Khoảng cách Mahalanobis bổ sung thông tin hiệp phương sai để chuẩn hóa độ lệch theo cấu trúc phân tán của không gian đặc trưng:
    \begin{equation}
    d_M(z,\mu_k)=\sqrt{(z-\mu_k)^\top(\Sigma+\lambda I)^{-1}(z-\mu_k)}
    \label{eq:mahalanobis}
    \end{equation}
    Trong đó, $z$ là vectơ biểu diễn của mẫu kiểm thử, $\mu_k$ là trung bình đặc trưng của lớp $k$ và $\Sigma$ là ma trận hiệp phương sai. Độ tin cậy của phép đo phụ thuộc vào chất lượng ước lượng hiệp phương sai và mức phù hợp của giả định phân phối.
\end{itemize}

\subsection{Học biểu diễn phục vụ phát hiện dữ liệu ngoài phân phối}
Sự thành công của các phương pháp dựa trên khoảng cách phụ thuộc rất lớn vào chất lượng của không gian đặc trưng. Nếu mô hình được huấn luyện bằng hàm entropy chéo thông thường, dữ liệu của các lớp không được tối ưu để co cụm lại, làm giảm hiệu quả của các phép đo khoảng cách. Để giải quyết vấn đề này, các kỹ thuật học biểu diễn được áp dụng.

\begin{itemize}
    \item \textbf{Học tương phản:} Học biểu diễn bằng cách kéo gần các mẫu tương đồng và đẩy xa các mẫu khác lớp trong không gian đặc trưng.
    \item \textbf{Mất mát tâm lớp \cite{CenterLoss}:} Giảm phương sai nội lớp bằng cách học một tâm biểu diễn cho mỗi lớp và ép các embedding tiến gần tâm tương ứng.
    \item \textbf{Mạng nguyên mẫu \cite{Snell2017ProtoNet}:} Phương pháp học theo hệ đo này biểu diễn mỗi lớp bằng một nguyên mẫu $c_k$:
    \begin{equation}
    c_k = \frac{1}{|S_k|}\sum_{x_i\in S_k}f_\theta(x_i)
    \label{eq:protonet}
    \end{equation}
    Trong quá trình huấn luyện, mạng được tối ưu hóa để cực tiểu hóa khoảng cách giữa đặc trưng của dữ liệu và nguyên mẫu tương ứng của nó.
\end{itemize}


\section{Các phương pháp giải thích và đánh giá độ trung thực}

Các phương pháp quy đóng gán mức ảnh hưởng cho từng thành phần đầu vào đối với một đầu ra cụ thể. Tích phân gradient \cite{Sundararajan2017IntegratedGradients} tích lũy gradient dọc theo đường nội suy từ một mốc tham chiếu đến đầu vào, nhờ đó có thể áp dụng cho vùng ảnh, đơn vị biểu diễn cục bộ và từ ngữ. Kết quả phản ánh mức phụ thuộc của đầu ra đối với mốc tham chiếu và lớp đích đã chọn, không mặc nhiên biểu thị quan hệ nhân quả hoặc vị trí tổn thương.

Độ trung thực của lời giải thích thường được kiểm tra bằng can thiệp đầu vào. Trong phép loại bỏ, các thành phần được xếp hạng cao được thay dần bằng mốc tham chiếu. Khi đó, xác suất lớp đích giảm nhanh hơn so với thứ tự ngẫu nhiên cung cấp bằng chứng thuận lợi hơn cho thứ tự quy đóng. Phép khôi phục thực hiện theo chiều ngược lại. Kết quả phụ thuộc vào mốc thay thế, tỷ lệ can thiệp và không gian biểu diễn, do đó cần báo cáo giao thức cùng đối chứng ngẫu nhiên thay vì chỉ dựa vào bản đồ nhiệt.

Trọng số chú ý mô tả cách mô-đun phân bổ tương tác trong một lần truyền thuận nhưng không mặc nhiên là lời giải thích trung thực cho quyết định. Ngoài độ trung thực, đánh giá trong y khoa còn cần phân biệt độ ổn định của bản đồ, mức phù hợp với chú thích tổn thương và nhận định của chuyên gia. Khóa luận vì vậy sử dụng tích phân gradient và can thiệp đầu vào để khảo sát nguồn thông tin liên hệ với dự đoán, đồng thời không diễn giải bản đồ quy đóng như bằng chứng định vị hoặc xác nhận lâm sàng.

\section{Bàn luận}

\subsection{Hạn chế của các nghiên cứu hiện nay}

Mặc dù đạt được nhiều kết quả tích cực, các nghiên cứu hiện có vẫn chưa giải quyết đồng thời các yêu cầu của bài toán phân lớp ảnh X-quang xương đa phương thức. Cùng với những hạn chế của mô hình BiomedCLIP khi phân lớp trực tiếp đã được trình bày trong Phần~\ref{subsubsec:BiomedCLIP_selection}, các vấn đề chính được tổng hợp như sau:

Thứ nhất, phần lớn mô hình y khoa đa phương thức được phát triển trên X-quang ngực hoặc các bộ dữ liệu lớn. Ảnh X-quang xương có phạm vi giải phẫu, tư thế chụp và kiểu tổn thương đa dạng, trong khi số lượng mẫu giữa các lớp thường mất cân bằng đáng kể. Vì vậy, hiệu quả trên các miền dữ liệu khác chưa thể được chuyển trực tiếp sang bài toán này.

Thứ hai, việc co toàn bộ ảnh về kích thước cố định có thể làm mất các đường gãy hoặc bất thường khu trú. Ngược lại, xử lý toàn bộ ảnh ở độ phân giải cao hoặc chia thành quá nhiều vùng làm số đơn vị biểu diễn và chi phí tính toán tăng nhanh. Do đó, cần kết hợp ngữ cảnh toàn cục với các vùng cục bộ có ích trong một ngân sách đơn vị biểu diễn được kiểm soát.

Thứ ba, nhiều nghiên cứu sử dụng báo cáo đọc ảnh hoặc kết luận chẩn đoán làm đầu vào văn bản. Tuy nhiên, các báo cáo này thường được hình thành sau khi bác sĩ hoặc bác sĩ chẩn đoán hình ảnh đã quan sát phim và có thể mô tả trực tiếp vị trí, hình thái hoặc mức độ tổn thương, chẳng hạn đường gãy, di lệch hay bất thường xương. Việc sử dụng các nguồn văn bản này trong giai đoạn suy luận có thể gây rò rỉ nhãn, khiến mô hình chủ yếu học cách ánh xạ một mô tả gần với kết luận chẩn đoán sang nhãn tương ứng. Trong bối cảnh triển khai thực tế, để tạo được báo cáo chi tiết, bác sĩ đã phải nhận diện tổn thương trước, do đó mô hình không còn đóng vai trò hỗ trợ phát hiện mà chủ yếu lặp lại thông tin đã được xác định. Ngược lại, các thông tin lâm sàng có sẵn trước chẩn đoán cuối cùng, như triệu chứng, tiền sử bệnh, vị trí đau hoặc cơ chế chấn thương, vẫn chưa được khai thác đầy đủ, mặc dù đây là những nguồn thông tin quan trọng giúp định hướng quá trình chẩn đoán. Bên cạnh đó, nhiều phương pháp chỉ nối hai vectơ đặc trưng toàn cục, nên chưa mô hình hóa được quan hệ cụ thể giữa các vùng ảnh và nội dung lâm sàng liên quan.

Thứ tư, mặc dù các mô hình nền tảng đạt hiệu quả cao, nhiều nghiên cứu vẫn lựa chọn tinh chỉnh toàn bộ tham số của mô hình. Cách tiếp cận này đòi hỏi chi phí tính toán lớn, đồng thời làm tăng nguy cơ quên tri thức đã được học trong giai đoạn tiền huấn luyện.

Cuối cùng, hầu hết các nghiên cứu chỉ tập trung tối ưu hiệu năng phân lớp mà chưa quan tâm đầy đủ đến khả năng phát hiện dữ liệu ngoài phân phối. Trong môi trường lâm sàng, mô hình có thể gặp các trường hợp bệnh lý hiếm, dữ liệu có chất lượng thấp hoặc các bất thường chưa xuất hiện trong tập huấn luyện. Nếu không có cơ chế phát hiện dữ liệu ngoài phân phối, mô hình vẫn có thể đưa ra dự đoán với độ tin cậy cao mặc dù kết quả không chính xác.

\subsection{Cơ sở lựa chọn và định hướng phương pháp đề xuất}
\label{subsec:ch2_method_direction}

Từ những hạn chế trên, khóa luận lựa chọn xây dựng hệ thống dựa trên việc kế thừa mô hình BiomedCLIP \cite{BiomedCLIP} kết hợp với thông tin lâm sàng để khai thác đồng thời đặc trưng hình ảnh và văn bản. Các lựa chọn thiết kế cụ thể được đưa ra dựa trên phân tích sau:

\textbf{Về xử lý ảnh độ phân giải cao:}
Việc co toàn bộ ảnh X-quang về kích thước đầu vào cố định có thể làm suy giảm các dấu hiệu tổn thương nhỏ, trong khi xử lý toàn bộ ảnh ở độ phân giải cao làm tăng đáng kể chi phí tính toán. Vì vậy, khóa luận lựa chọn chiến lược toàn cục-cục bộ, trong đó ảnh toàn cục duy trì bố cục giải phẫu chung, còn bốn vùng cục bộ được lựa chọn nhằm bảo toàn các vùng giàu thông tin. Các đặc trưng cục bộ được bổ sung thông tin vị trí và nén về số lượng đơn vị biểu diễn cố định để cân bằng giữa khả năng giữ lại chi tiết và chi phí xử lý.

\textbf{Về thích nghi mô hình nền tảng:}
LoRA \cite{Hu2022LoRA} được sử dụng để thích nghi hai bộ mã hóa của BiomedCLIP với miền ảnh X-quang xương mà không cần cập nhật toàn bộ tham số nền. Lựa chọn này giúp giảm chi phí huấn luyện, hạn chế nguy cơ quá khớp trên các lớp có ít mẫu và bảo toàn phần lớn tri thức y sinh đã được học trong giai đoạn tiền huấn luyện.

Bên cạnh đó, hàm mất mát tương phản InfoNCE nguyên bản xem các mẫu khác nhau trong cùng lô là các cặp âm, kể cả khi chúng thuộc cùng một lớp hoặc có nội dung lâm sàng tương đồng. Do đó, khóa luận sử dụng hàm mất mát khớp ngữ nghĩa kế thừa từ MedCLIP và xây dựng phân phối đích mềm theo nhãn lớp. Mục tiêu này vẫn ưu tiên sự tương hợp của cặp ảnh-văn bản tương ứng, đồng thời giảm ảnh hưởng của các cặp âm tính giả giữa những mẫu có quan hệ ngữ nghĩa gần nhau.

\textbf{Về khai thác và dung hợp thông tin lâm sàng:}
Thiết lập phân lớp không mẫu của BiomedCLIP chỉ đối sánh ảnh với các câu nhắc đại diện cho lớp và không sử dụng thông tin riêng của từng người bệnh. Phương pháp đề xuất bổ sung nhánh văn bản để mã hóa những thông tin lâm sàng có sẵn trước kết luận chẩn đoán, chẳng hạn lý do vào viện, triệu chứng và bệnh sử. Báo cáo đọc ảnh, chẩn đoán xác định và nhãn phân lớp không được đưa trực tiếp vào đầu vào nhằm hạn chế rò rỉ nhãn.

Thay vì chỉ ghép nối hai biểu diễn toàn cục, khóa luận lựa chọn cơ chế chú ý chéo hai chiều để đặc trưng ảnh lựa chọn các thông tin lâm sàng liên quan, đồng thời đặc trưng văn bản định hướng mô hình đến các vùng ảnh phù hợp. Thiết kế này cho phép mô hình hóa quan hệ cụ thể giữa dấu hiệu X-quang và ngữ cảnh lâm sàng.

\textbf{Về phân lớp đa lớp và mất cân bằng dữ liệu:}
Phương pháp sử dụng đầu phân lớp dựa trên tâm lớp thực nghiệm, được phát triển từ ý tưởng của mạng nguyên mẫu, trong đó mỗi lớp được đại diện bằng trung bình các vectơ biểu diễn thuộc lớp đó thay vì một nguyên mẫu được học tự do. Độ tương đồng cosine với từng tâm được chuyển thành logit và kết hợp với độ lệch học được theo lớp trước khi áp dụng Softmax. Vì vậy, nhãn dự đoán không nhất thiết luôn là tâm có độ tương đồng lớn nhất. Hàm mất mát entropy chéo có trọng số được sử dụng nhằm giảm ảnh hưởng của sự chênh lệch số lượng mẫu giữa các lớp.

\textbf{Về phát hiện dữ liệu ngoài phân phối:}
Ngoài nhiệm vụ phân lớp, khóa luận bổ sung cơ chế phát hiện dữ liệu ngoài phân phối ở giai đoạn hậu xử lý mà không yêu cầu mẫu ngoài phân phối trong quá trình huấn luyện, hoạt động dựa trên cơ chế phân loại trong không gian đặc trưng được xây dựng từ bộ phân loại. Khoảng cách Mahalanobis được tính giữa biểu diễn của mẫu mới và các trung bình lớp thực nghiệm, có xét đến cấu trúc hiệp phương sai của không gian đặc trưng. Mẫu có điểm vượt ngưỡng được gắn thêm cảnh báo ngoài phân phối. Trong phạm vi khóa luận, cảnh báo này không thay đổi nhãn phân lớp nội tại và chưa được sử dụng để tự động từ chối dự đoán. Ngưỡng được hiệu chỉnh trên tập xác thực.

\textbf{Về quy trình huấn luyện:}
Từ các lựa chọn trên, hệ thống được tổ chức thành hai giai đoạn ngoại tuyến:

\begin{itemize}
    \item \textbf{Giai đoạn 1 - Thích nghi và căn chỉnh biểu diễn:}
    hai bộ mã hóa của BiomedCLIP được thích nghi thông qua LoRA bằng hàm mất mát khớp ngữ nghĩa. Mô-đun dung hợp và đầu phân lớp chưa được tối ưu trong giai đoạn này.

    \item \textbf{Giai đoạn 2 - Dung hợp và phân lớp:}
    các bộ mã hóa tiếp tục được thích nghi thông qua LoRA, đồng thời mô-đun xử lý đặc trưng cục bộ và cơ chế chú ý chéo hai chiều được tối ưu cho nhiệm vụ phân lớp. Các tâm lớp thực nghiệm được cập nhật từ đặc trưng của tập huấn luyện và được sử dụng cho cả phân lớp lẫn phát hiện dữ liệu ngoài phân phối.
\end{itemize}

Các định hướng trên được hiện thực hóa trong phương pháp XBone-Net. Kiến trúc chi tiết, quy trình huấn luyện và cơ chế suy luận được trình bày trong Chương~\ref{Chapter3}.
